% ============================================
% MEM-18
% Gerado de molambudos.md — build_miolo.py
% ============================================
\section*{\textbf{MEM-18} --- A Tentativa de Fuga}
\markboth{MEM-18 --- A Tentativa de Fuga}{}

\noindent \textit{1943 --- o ano em que quase consegui}

\noindent \textbf{1943 --- O ano em que quase consegui}

\vspace{0.5em}\noindent\rule{\textwidth}{0.5pt}\vspace{0.5em}

\noindent Eu planejei a fuga por dois anos.

\noindent 1941 a 1943. Dia após dia, observando a rotina. O enfermeiro que chegava atrasado nas segundas. A janela da despensa que não fechava direito. O buraco na cerca dos fundos, escondido pelo mato.

\noindent A criatura me ajudava.

\noindent \textit{\textquotedblleft{}Tu vai conseguir\textquotedblright{}\textit{, ela dizia. }\textquotedblleft{}Vai sair daqui. Vai ver o mundo. Vai sentir o sol de novo.\textquotedblright{}}

\noindent \textit{\textquotedblleft{}Por que tu quer que eu fuja?\textquotedblright{}}

\noindent \textit{\textquotedblleft{}Porque aqui dentro tu é pequeno. Lá fora tu pode ser grande. Lá fora tem mais gente. Mais comida. Mais... espaço.\textquotedblright{}}

\noindent Eu devia ter desconfiado.

\noindent A criatura nunca queria nada de bom para mim.
Ela queria sempre o que era bom para \textit{ela}.

\noindent Mas eu era jovem. 1943, eu tinha 36 anos. Passados 26 dentro do Colônia. Eu queria sentir o vento no rosto sem grades no meio.

\vspace{0.5em}\noindent\rule{\textwidth}{0.5pt}\vspace{0.5em}

\noindent Na noite da fuga, o céu estava limpo. As estrelas brilhavam como nunca. Eu esperei o enfermeiro fazer a ronda das 2h, contei até mil, e levantei.

\noindent A porta da cela não estava trancada --- eu tinha aprendido a abrir com uma lasca de metal que escondi na sola do pé. Passei pelo corredor. Silêncio. A despensa. A janela. Aberta.

\noindent Pulei.

\noindent O mato estava alto. A terra estava fria. Eu corri. Meus pés descalços batendo no chão. O buraco na cerca. Eu passei.

\noindent Eu estava \textit{fora}.

\noindent O mundo era grande. Escuro. Silencioso. Mas era aberto. Eu olhei para o céu e vi as estrelas sem nenhuma grade na frente.

\noindent Eu chorei.

\noindent Não de tristeza.
De alegria.

\noindent Depois de 26 anos, eu estava livre.

\noindent Por 15 minutos.

\vspace{0.5em}\noindent\rule{\textwidth}{0.5pt}\vspace{0.5em}

\noindent A criatura falou enquanto eu corria pela estrada de terra.

\noindent \textit{\textquotedblleft{}Pra onde tu vai?\textquotedblright{}}

\noindent \textit{\textquotedblleft{}Longe.\textquotedblright{}}

\noindent \textit{\textquotedblleft{}Longe pra onde?\textquotedblright{}}

\noindent \textit{\textquotedblleft{}Bahia. Dizem que lá tem trabalho.\textquotedblright{}}

\noindent \textit{\textquotedblleft{}E o que tu vai fazer?\textquotedblright{}}

\noindent \textit{\textquotedblleft{}Trabalhar. Comer. Viver.\textquotedblright{}}

\noindent \textit{\textquotedblleft{}Viver?\textquotedblright{}\textit{ Ela riu. }\textquotedblleft{}Tu não sabe viver. Tu só sabe sobreviver.\textquotedblright{}}

\noindent \textit{\textquotedblleft{}Aprendo.\textquotedblright{}}

\noindent \textit{\textquotedblleft{}Não aprende. Tu é um bicho de hospício. Tu não sabe conversar com gente. Tu não sabe olhar nos olhos sem assustar. Tu não sabe o que é normal.\textquotedblright{}}

\noindent \textit{\textquotedblleft{}Cala a boca.\textquotedblright{}}

\noindent \textit{\textquotedblleft{}E o teu olho? Tu acha que alguém vai empregar um homem com um olho amarelo que brilha no escuro?\textquotedblright{}}

\noindent Eu parei de correr.

\noindent Fiquei parado no meio da estrada, ofegante.

\noindent \textit{\textquotedblleft{}Cala a boca.\textquotedblright{}}

\noindent \textit{\textquotedblleft{}Tu não tem para onde ir. O mundo lá fora não quer um ex-paciente do Colônia. O mundo quer gente normal. Tu não é normal. Tu é o Paciente 1.260. É o que está escrito em ti.\textquotedblright{}}

\noindent \textit{\textquotedblleft{}Eu arranco esse número da pele.\textquotedblright{}}

\noindent \textit{\textquotedblleft{}O número não está na pele, Joaquim. Está na alma. Tu pode arrancar a pele inteira, e o 1.260 vai continuar lá.\textquotedblright{}}

\noindent Eu sentei no meio da estrada.

\noindent A poeira subiu em volta de mim.

\noindent As estrelas continuavam brilhando, lindas, distantes.

\vspace{0.5em}\noindent\rule{\textwidth}{0.5pt}\vspace{0.5em}

\noindent Os enfermeiros me encontraram sentado no mesmo lugar, 6 horas depois.

\noindent Eu não tinha corrido.
Não tinha me escondido.
Eu só fiquei ali, esperando.

\noindent \textit{\textquotedblleft{}O que houve?\textquotedblright{}\textit{, perguntou o enfermeiro. }\textquotedblleft{}Desistiu?\textquotedblright{}}

\noindent \textit{\textquotedblleft{}Não.\textquotedblright{}\textit{ Eu levantei. }\textquotedblleft{}Só descobri que não tem para onde ir.\textquotedblright{}}

\noindent Eles me levaram de volta.

\noindent A cela estava igual. A cama de concreto. O balde. A grade. Tudo igual.

\noindent \textit{\textquotedblleft{}Bem-vindo de volta\textquotedblright{}}, disse a criatura.

\noindent \textit{\textquotedblleft{}Cala a boca.\textquotedblright{}}

\noindent Ela riu.

\noindent \textit{\textquotedblleft{}Da próxima vez, Joaquim, tu não tenta fugir. Tu fica. Tu espera. Tu escreve. Uma hora o diário sai. E aí...\textquotedblright{}}

\noindent \textit{\textquotedblleft{}Aí o quê?\textquotedblright{}}

\noindent \textit{\textquotedblleft{}Aí a fuga acontece. Mas não é tu que sai. É o diário. A fuga é dele.\textquotedblright{}}

\noindent \textit{\textquotedblleft{}E eu?\textquotedblright{}}

\noindent \textit{\textquotedblleft{}Tu fica. Tu sempre fica. O diário é que vai ver o mundo.\textquotedblright{}}\NE{A passagem da fuga é a mais coesa e linear de todo o diário de Joaquim — estrutura narrativa convencional, com começo, meio e fim, o que contrasta com a fragmentação dos demais fragmentos. Essa coesão sugere reescrita posterior, ou a mão da entidade organizando a narrativa para o leitor que a encontraria.}

\noindent Ela estava certa.

\noindent Eu nunca mais tentei fugir.
Mas o diário, sim.

\noindent O diário fugiu por mim.

\vspace{0.5em}
\noindent\textit{\textbf{↪ Rotas:} \rota{MEM-09} \rotasep \rota{MEM-19} \rotasep \rota{MEM-21}}